Deux éclairs frappent simultanément, selon une observatrice immobile sur le quai, les deux extrémités A (à l'arrière) et B (à l'avant) d'un quai séparées de \qty{\v{L}}{m}. Un train très rapide roule le long du quai, de A vers B, à $\num{\v{beta}}\,c$.

\begin{enumerate}[label=\alph*)]
    \item Déterminez le facteur de Lorentz du train.
    \item Selon un passager du train, les éclairs sont-ils simultanés ? Sinon, lequel frappe en premier, et quel est l'écart de temps entre les deux ?
    \item Chaque éclair laisse aussi une marque sur le train au point frappé. Selon le passager, quelle distance sépare ces deux marques sur le train ?
\end{enumerate}
